\documentclass[11pt]{article}

\usepackage[margin=1in]{geometry}
\usepackage{amsmath,amssymb}
\usepackage{booktabs}
\usepackage{enumitem}
\usepackage{longtable}
\usepackage{array}
\usepackage[T1]{fontenc}
\usepackage[utf8]{inputenc}
\usepackage{xcolor}
\usepackage{hyperref}
\usepackage{url}

\hypersetup{
  pdftitle={Implementation-Grade Survey for Hypothesis-Tree Derivation Search in MathDevMCP},
  hidelinks
}

\newcommand{\MathDevMCP}{\textsc{MathDevMCP}}
\newcommand{\Lean}{\textsc{Lean}}
\newcommand{\Coq}{\textsc{Coq}}
\newcommand{\Isabelle}{\textsc{Isabelle}}
\newcommand{\Sage}{\textsc{SageMath}}
\newcommand{\SymPy}{\textsc{SymPy}}
\newcommand{\MCTS}{\textsc{MCTS}}
\newcommand{\UCT}{\textsc{UCT}}
\newcommand{\LLM}{\textsc{LLM}}
\newcommand{\API}{\textsc{API}}
\setlength{\emergencystretch}{3em}

\title{Implementation-Grade Survey for Hypothesis-Tree Derivation Search in \MathDevMCP}
\author{}
\date{2026-07-07}

\begin{document}
\maketitle

\begin{abstract}
\MathDevMCP{} needs a stronger engine for hard mathematical document repair.
The current workflows can locate gaps, discover some missing assumptions, and
abstain safely, but difficult repairs often require a search over alternative
formalizations, assumption sets, derivation decompositions, and backend routes.
This note surveys the relevant theorem-proving and language-agent literature,
inspects official code where available, and turns the findings into an
implementation decision document. The recommendation is to build a
\emph{MathDevMCP-native evidence-grounded hypothesis tree}: use deterministic
source parsers and backends as the environment, borrow best-first/MCTS-style
branching and proof-state data structures from existing systems, and keep
Lean/Pantograph/LeanDojo integrations optional adapters rather than required
core dependencies.
\end{abstract}

\tableofcontents
\newpage

\section{Purpose, Evidence Contract, and Non-Claims}

\paragraph{Question.}
Which existing algorithms and codebases should guide or support a
\MathDevMCP{} feature that searches over derivation hypotheses and
missing-assumption branches, records blocker nodes, and delegates checkable
steps to deterministic backends?

\paragraph{Why this matters.}
The weak report failure is not only a rendering problem. A report line such as
``justify the reconstructed equality from context'' is weak because the system
has not searched for the missing derivation. For hard documents, the missing
piece may be an assumption, a definition, a source-span reconstruction, a
rewrite route, a measure-theoretic condition, a shape/type condition, or a
formalization choice. A single linear attempt is structurally underpowered.

\paragraph{Primary criterion.}
This document should be self-contained enough for an implementation decision.
For each major candidate it records:

\begin{itemize}[leftmargin=*]
\item the problem the paper/tool solves;
\item the actual state/action/search/scoring algorithm;
\item what code exists and what was inspected locally;
\item whether the code appears directly installable or usable;
\item whether \MathDevMCP{} should reuse it, integrate against it, adapt the
  algorithm only, or avoid it.
\end{itemize}

\paragraph{Veto diagnostics.}
This note must not treat \LLM{} output as proof, must not claim installability
without a local check, must not recommend direct reuse without inspecting
dependencies and runtime shape, and must not confuse proof-assistant theorem
proving with mathematical-document repair.

\paragraph{Not concluded.}
This is not a benchmark. It does not prove that any external code can solve the
credit-card NPV or risky-debt documents. It is a decision artifact for the next
implementation plan.

\section{Local Baseline and the Actual Problem}

\subsection{Current \MathDevMCP{} Behavior}

Local inspection shows these relevant workflows:

\begin{itemize}[leftmargin=*]
\item \texttt{derive\_or\_refute}: split target into lhs/rhs, discover
  assumptions, route to a backend, optionally run counterexample search, and
  return one result.
\item \texttt{derive\_from}: wrap the low-level derivation/refutation result
  into a high-level result with gaps, proposals, tool uses, and non-claims.
\item \texttt{assumptions\_for}: propose route-required assumptions and
  possible sufficient assumption sets.
\item \texttt{audit\_and\_propose\_fix}: audit labels and use
  \texttt{ProcessPoolExecutor} across labels.
\item \texttt{audit\_math\_document\_rigor}: render a document-level report
  with concrete-repair and diagnostic-abstention ledgers.
\end{itemize}

The repo has parallelism across labels. It does not yet have a within-problem
search tree that expands alternative assumption branches, subgoals, and backend
routes for one hard derivation.

\subsection{The Mathematical Search Problem}

For a hard document target, the system receives a claim \(T\), source context
\(C\), and known assumptions \(A_0\). It wants to produce one of:

\begin{enumerate}[leftmargin=*]
\item a derivation of \(T\) from source-backed assumptions and checked steps;
\item a refutation/counterexample;
\item a best-progress branch with exact remaining blockers;
\item a budget-exhausted search tree that is still useful to a human or agent.
\end{enumerate}

The search state is not simply a Lean proof state. For \MathDevMCP{} it should
include:

\[
  s = (T, C, A, G, D, B, E, P)
\]

where \(A\) is the current assumption set, \(G\) is the list of subgoals,
\(D\) is the partial derivation, \(B\) is the blocker ledger, \(E\) is backend
evidence, and \(P\) is the candidate patch/source-span state.

Actions include:

\begin{itemize}[leftmargin=*]
\item add a candidate assumption set;
\item split the target into subgoals;
\item recover or expand a source span;
\item rewrite an expression using a definition or accounting identity;
\item run a backend check;
\item ask a proof assistant for tactic/proof-state feedback;
\item mark a branch blocked with evidence.
\end{itemize}

The success condition is not ``the agent sounds confident.'' It is:

\begin{quote}
Every step on a winning branch is either backend checked, source-backed and
accepted as an explicit assumption, or explicitly marked as an unchecked
modeling convention.
\end{quote}

\subsection{Why the Current Reports Regress}

The weak-report regression has a structural cause. The current high-level
reports can say that a problem is unresolved, and they can sometimes produce a
local assumption proposal, but they do not yet maintain a search state across
multiple attempted repairs for the same target. When a branch fails, the
failure is summarized rather than converted into a new, typed subproblem. This
causes three concrete pathologies:

\begin{enumerate}[leftmargin=*]
\item \textbf{One-route collapse.} A target is sent through one derivation
  route. If that route is inconclusive, the report falls back to generic
  advice even when another route is available.
\item \textbf{Missing blocker memory.} A failure such as ``conditional law not
  defined'' is not retained as a node that can be expanded by postulating a
  finite-state kernel, a measurable transition kernel, or a deterministic
  sequence convention.
\item \textbf{No branch-level evidence accounting.} The renderer cannot say
  which assumptions made which derivation step work, because the underlying
  result is not a branch containing assumptions, checked steps, rejected
  alternatives, and remaining blockers.
\end{enumerate}

The tool that is missing is therefore not only ``a better renderer.'' The
missing tool is a reusable search controller whose output is already shaped as
an agent-consumable repair report.

\subsection{Minimum Agent-Consumable Output}

For a hard target, the function should produce a result with this contract:

\begin{verbatim}
{
  "question": "can we derive target T from context C?",
  "budget": {"max_nodes": ..., "max_depth": ..., "timeouts": ...},
  "best_branch": {
    "status": "proved|refuted|partial|blocked|budget_exhausted",
    "assumptions_added": [...],
    "derivation_steps": [
      {
        "claim": "...",
        "justification": "...",
        "checker": "parser|sympy|sage|lean|source_assumption",
        "checker_status": "certified|accepted_assumption|not_certified",
        "evidence_ref": "..."
      }
    ],
    "patch": {"assumption_latex": "...", "replacement_latex": "..."}
  },
  "blocked_branches": [...],
  "rejected_branches": [...],
  "non_claims": [...]
}
\end{verbatim}

The result should be directly consumable by an agent: it should say where the
claim is, what was tried, which assumptions were introduced, why those
assumptions close a mathematical gap, what backend checked which step, and what
precise patch should be made if the branch is selected.

\subsection{Example: Finite-State Conditional Expectation Branch}

A typical stochastic-economics repair should not say only ``assume
integrability.'' It should show the route. Suppose a document uses
\[
  E_t[m_{t+1} X_{t+1}]
\]
without defining the conditional law. One branch can postulate a finite-state
shock space \(\Omega=\{\omega_1,\ldots,\omega_K\}\) with conditional
probabilities \(p_k(s_t)\). That branch transforms the expectation into
\[
  E_t[m_{t+1}X_{t+1}]
  =
  \sum_{k=1}^K p_k(s_t)\,
  m(s_t,\omega_k) X(s_t,\omega_k).
\]
The branch then has concrete obligations:

\begin{enumerate}[leftmargin=*]
\item \(p_k(s_t)\ge 0\) and \(\sum_k p_k(s_t)=1\), so the conditional law is
  defined.
\item \(m(s_t,\omega_k)X(s_t,\omega_k)\) is finite for every \(k\), so the
  finite expectation is integrable.
\item The algebraic rewrite from expectation notation to weighted finite sum is
  a definition expansion, not an economic theorem.
\item Any subsequent equality can be checked by \SymPy/\Sage{} as an algebraic
  identity over the finite sum, or blocked with the exact missing identity.
\end{enumerate}

This is the level of repair detail the new search controller must make
available to the report layer.

\section{Code Inspection Protocol}

All cloned code lives under:

\begin{verbatim}
.localresources/hypothesis_search_survey/code/
\end{verbatim}

The inspection used:

\begin{itemize}[leftmargin=*]
\item \texttt{git ls-remote} to check public repository availability;
\item shallow \texttt{git clone --depth 1} for official repositories;
\item README/license/dependency/API inspection;
\item lightweight local probes for installed Python modules and executables;
\item no heavy model download, Coq/Isabelle install, Hugging Face download, or
  broad environment mutation.
\end{itemize}

Local environment probes found:

\begin{itemize}[leftmargin=*]
\item \Sage{} is available at \texttt{/usr/bin/sage}.
\item \SymPy{} imports in active Python.
\item \Lean/\texttt{lake}/\texttt{elan} binaries exist, but local version
  checks timed out in the active diagnostic path.
\item \texttt{lean\_dojo}, \texttt{pantograph}, \texttt{tot}, \texttt{torch},
  \texttt{transformers}, and \texttt{openai} were not importable in the active
  Python used by the probe.
\item Existing \MathDevMCP{} code already has LeanDojo readiness/policy
  wrappers, including the requirement that LeanDojo output must be reconstructed
  and checked by direct Lean before certification.
\end{itemize}

\section{Code and Paper Evidence Ledger}

\subsection{Local Paper Copies}

The following papers were downloaded locally because they materially affect the
implementation decision:

\begin{itemize}[leftmargin=*]
\item \path{.localresources/hypothesis_search_survey/papers/leandojo-2306.15626.pdf}.
\item \path{.localresources/hypothesis_search_survey/papers/pantograph-2410.16429.pdf}.
\item \path{.localresources/hypothesis_search_survey/papers/htps-2205.11491.pdf}.
\item \path{.localresources/hypothesis_search_survey/papers/tree-of-thoughts-2305.10601.pdf}.
\item \path{.localresources/hypothesis_search_survey/papers/deepseek-prover-v15-2408.08152.pdf}.
\end{itemize}

\subsection{Clone Manifest}

\begin{center}
\small
\begin{tabular}{lll}
\toprule
Repository & Local HEAD & Primary inspection target \\
\midrule
LeanDojo & \texttt{7a9f600b250a} & README, package metadata \\
LeanDojo-v2 & \texttt{936ea0dd32ff} & README, \texttt{pyproject.toml}, prover tree \\
ReProver & \texttt{fd6d99c01e8b} & \texttt{prover/search\_tree.py}, \texttt{proof\_search.py} \\
PyPantograph & \texttt{f8aee320ee55} & \texttt{pantograph/server.py}, \texttt{search.py} \\
aesop & \texttt{24fa6b359974} & README, \texttt{lakefile.toml}, \texttt{lean-toolchain} \\
LeanCopilot & \texttt{2458f7339df4} & README, \texttt{lakefile.lean} \\
CoqGym & \texttt{a739d99cdf5b} & README, setup/dependency structure \\
Evariste & \texttt{ef43511f2d23} & README, setup files, license \\
tree-of-thought-llm & \texttt{8050e67d0e3a} & \texttt{src/tot/methods/bfs.py}, requirements \\
LanguageAgentTreeSearch & \texttt{853d81614607} & \texttt{programming/mcts.py}, requirements \\
DeepSeek-Prover-V1.5 & \texttt{2c4ba9119eef} & README, \texttt{configs/RMaxTS.py}, requirements \\
lean-explore & \texttt{3a52d6b9922f} & README, MCP/API docs, search models \\
LeanSearch-v2 & \texttt{94f4888cbaf9} & README, client, retrieval pipeline \\
jixia & \texttt{755fde27a9cf} & README, Lean static-analysis plugins \\
\bottomrule
\end{tabular}
\end{center}

\subsection{Local Probes}

Bounded probes were used to avoid mutating the active environment:

\begin{itemize}[leftmargin=*]
\item Active Python can import \SymPy{} but cannot import
  \texttt{lean\_dojo}, \texttt{lean\_dojo\_v2}, \texttt{pantograph},
  \texttt{tot}, \texttt{torch}, \texttt{transformers}, \texttt{openai}, or
  \texttt{ray}.
\item \texttt{python3 -m py\_compile} passed for the inspected ReProver
  search files, PyPantograph search/server files, ToT BFS file, LATS
  \texttt{programming/mcts.py}, and DeepSeek-Prover-V1.5 launch/config files.
\item \texttt{python3 -m compileall} passed for the inspected
  \texttt{lean-explore}, \texttt{LeanSearch-v2}, and \texttt{jixia} source
  trees. This is a syntax/readiness probe only; it is not an install or runtime
  proof.
\item A broad compile of the LATS programming directory failed on vendored
  \texttt{human\_eval/execution.py}; therefore only the MCTS controller file is
  treated as locally syntax-checked.
\item \texttt{python3 -m mathdevmcp.cli doctor} found \Sage{} and \SymPy{} but
  reported Lean readiness as inconclusive in the active diagnostic path.
\item No \texttt{pip install}, \texttt{lake build}, model download, Docker
  build, OPAM/Coq setup, or GPU run was attempted.
\end{itemize}

\subsection{Readiness Classification}

This survey uses three implementation classifications:

\begin{description}[style=nextline,leftmargin=0pt]
\item[Use directly.] The code is small enough, locally installable or already
usable, and matches the \MathDevMCP{} contract without forcing heavy model or
proof-assistant setup.
\item[Integrate/test against.] The code exposes a useful runtime boundary, but
requires an environment gate, optional backend env, model/API token, or proof
assistant project before it can be used safely.
\item[Not recommended as dependency.] The code is too heavy, stale, non-running,
not proof-specific, license-limited, or mismatched with document-repair. Its
algorithm may still be worth adapting.
\end{description}

\subsection{Candidate-by-Candidate Result}

\begin{description}[style=nextline,leftmargin=0pt]
\item[LeanDojo original.] \textbf{Classification: not recommended as a new
dependency.} The README says the original package is deprecated and directs new
projects to LeanDojo-v2. It remains important for existing compatibility and for
understanding the ReProver paper, but new work should not bind the core search
controller to it.

\item[LeanDojo-v2.] \textbf{Classification: integrate/test against later.}
It is a complete framework for tracing, training, retrieval, agents, and
Pantograph-based proving, but its package metadata requires Python 3.11 and
heavy dependencies including PyTorch, Lightning, Transformers, DeepSpeed, TRL,
PEFT, Ray, OpenAI, PyGithub, and datasets. It belongs in an isolated backend
environment, not the core \MathDevMCP{} import path.

\item[ReProver.] \textbf{Classification: adapt algorithm only now; integrate
later only through a LeanDojo lane.} The local files are valuable because they
show a compact proof-search skeleton: status enum, internal nodes, edges,
cumulative priority, best-first queue, tactic execution, status propagation,
and proof extraction. Direct reuse pulls in LeanDojo, PyTorch, Ray, vLLM, and
model checkpoints, so the immediate value is architectural.

\item[PyPantograph.] \textbf{Classification: integrate/test against after an
environment gate.} It is the best optional proof-state adapter candidate. Its
Python package has no ordinary Python dependencies, but the build script
bundles a Lean REPL and requires matching Lean/lake setup. It should implement
a \texttt{ProofStateBackend} after direct Lean readiness is stable.

\item[Aesop.] \textbf{Classification: integrate/test against as a Lean-side
tactic after a Lean project gate; adapt algorithm now.} Aesop is not a Python
package. Its safe/unsafe rule discipline and best-first AND/OR tree are a
strong template for \MathDevMCP{} expansion operators.

\item[Lean Copilot.] \textbf{Classification: not core; possible later
experiment.} It is useful as a Lean-side tactic/model interface reference, but
requires Lean project setup, C++/CMake machinery, model assets, and Git LFS.

\item[CoqGym.] \textbf{Classification: not recommended as dependency.} It is
useful as a replayable proof-state dataset reference. Direct use requires a
Coq/SerAPI/CoqHammer/OPAM/Docker stack that does not match the current Lean/Sage
lane.

\item[Evariste / HTPS.] \textbf{Classification: not recommended as dependency.}
The README says the code does not run out of the box after Meta-internal
components were removed, and the majority license is CC-BY-NC. The AND/OR
hypergraph algorithm is still important.

\item[Tree of Thoughts.] \textbf{Classification: not recommended as dependency.}
The code is simple and the BFS file syntax-checks, but the system is prompt/API
centric and depends on the old OpenAI package. Borrow only the
generate/evaluate/select abstraction.

\item[Language Agent Tree Search.] \textbf{Classification: not recommended as
dependency.} The MCTS file syntax-checks and is a good control-loop reference,
but the broader programming tree has a syntax error in vendored HumanEval code
and the project depends on API/task executors. Borrow the feedback/backprop
ledger idea.

\item[DeepSeek-Prover-V1.5.] \textbf{Classification: not core; strategic
reference and possible future high-budget benchmark.} The repo is public and
the inspected config exposes RMaxTS with large settings: batch size 512, Lean
concurrency 64, sample count 6400, and 256 search processes. Requirements
include PyTorch, Transformers, vLLM, FlashAttention, and GPU-oriented inference.
It is evidence that hard proof search benefits from verifier feedback and many
attempts, not a drop-in document-repair component.

\item[HunyuanProver.] \textbf{Classification: paper/strategy reference only for
this survey.} A GitHub repository search for \texttt{HunyuanProver} returned no
obvious official repo in the bounded probe. Do not plan an integration until an
official codebase is identified and inspected.

\item[LeanExplore.] \textbf{Classification: readily borrowable integration
pattern; possible optional dependency after environment review.} The cloned
repo is a Python package with Apache-2.0 license, a hosted API client, a local
hybrid retrieval backend, and an MCP server. Its MCP design is especially
relevant: agents first call a low-token \texttt{search\_summary} tool and then
fetch only the source, docstring, module, dependencies, or description for the
few declarations they need \cite{asher2025leanexplore}. This is directly
portable to \MathDevMCP{}.

\item[LeanSearch-v2.] \textbf{Classification: borrow algorithms and client
shape now; do not make the local server a core dependency.} The repo targets
global premise retrieval with a standard retrieval endpoint and a reasoning
mode. Its local standard pipeline requires GPU-backed cuVS, Qwen embedding and
reranking models, and Python 3.11; its thin HTTP client is lightweight. The
most useful immediate ideas are global premise retrieval, decompose/retrieve/
filter/judge loops, and the rule that proof attempts must use only retrieved
tools/theorems \cite{gao2026leansearchv2}.

\item[jixia.] \textbf{Classification: borrow source-analysis contract now;
optional Lean adapter later.} jixia is a Lean 4 static-analysis tool that emits
JSON for declarations, symbols, elaboration information, proof-state lines, and
ASTs. Its README stresses exact Lean-version matching. This makes it a strong
template for source-range and dependency ledgers, but not a zero-setup core
dependency \cite{jixiagithub}.
\end{description}

\section{Second-Pass Citation Snowball and Borrowability Audit}

\subsection{Scope and Limitation}

After the first survey pass, a citation-snowball script was written under
\path{.localresources/hypothesis_search_survey/snowball/} to query Semantic
Scholar for references and citations of the seed papers. The API returned HTTP
429 for all seed papers in the bounded run, and
\path{candidate_papers.json} is therefore empty. This note must not claim a
complete forward/backward citation graph.

The fallback pass used three evidence sources:

\begin{enumerate}[leftmargin=*]
\item local PDF text mining from LeanDojo/ReProver, Pantograph, HTPS,
  Tree-of-Thoughts, and DeepSeek-Prover-V1.5;
\item targeted arXiv/GitHub searches for systems named in those related-work
  and comparison sections;
\item shallow local clone inspection for the most directly borrowable code
  surfaces.
\end{enumerate}

This is a bounded implementation survey, not an exhaustive bibliometric review.
Its purpose is to decide which mechanisms \MathDevMCP{} can borrow soon.

\subsection{What the Snowball Pass Added}

The first pass already covered major proof-search families, but it underweighted
premise retrieval and declaration search. The second pass surfaced four
borrowable clusters:

\begin{enumerate}[leftmargin=*]
\item \textbf{Agent-token-efficient declaration search.} LeanExplore exposes a
  low-token MCP workflow: search summaries first, then per-field source,
  docstring, dependency, or module fetches. This matches \MathDevMCP{}'s goal of
  reducing agent hallucination and token use.
\item \textbf{Global premise retrieval.} LeanSearch-v2 and ReProver both show
  that hard proof search needs a premise index before generation. LeanSearch-v2
  is too heavy as a core dependency, but its standard-client boundary and
  reasoning retrieval loop are reusable design patterns.
\item \textbf{Source-range and dependency extraction.} jixia, LeanDojo, and
  CoqGym all treat source ranges, dependencies, proof states, and replay logs as
  first-class data. \MathDevMCP{} should do the same for LaTeX labels,
  assumptions, equations, and derived obligations.
\item \textbf{Partial-progress proof search.} DeepSeek-Prover-V1.5's
  truncate-and-resume search and RMaxTS exploration mechanism reinforce the
  blocker-node design: a failed attempt should become a reusable partial path,
  not a generic report sentence.
\end{enumerate}

\subsection{Coverage Matrix for Readily Borrowable Mechanisms}

\begin{small}
\begin{longtable}{p{0.18\linewidth}p{0.26\linewidth}p{0.23\linewidth}p{0.23\linewidth}}
\toprule
Family or system & Mechanism & Borrowability for \MathDevMCP{} & Concrete action \\
\midrule
\endhead
LeanExplore & MCP/API search with \texttt{search\_summary} plus per-field
getters for source, dependencies, docstrings, descriptions, and modules &
High. The interaction pattern is directly agent-consumable and token-frugal;
local backend is optional & Add a summary-first context-search provider with
field-fetch methods; later add optional LeanExplore adapter \\
\addlinespace
LeanSearch-v2 & Global premise retrieval; standard HTTP client; reasoning mode
that decomposes, retrieves, filters, and judges candidate premises & Medium.
The client shape and reasoning loop are useful; the full local server is GPU
and model heavy & Borrow the decompose--retrieve--filter--judge contract and
``use only retrieved tools'' rule; keep GPU server outside core \\
\addlinespace
jixia & Non-intrusive Lean source analysis: declarations, symbols, elaboration,
line proof states, AST dumps & Medium. Great contract for exact source
provenance; requires exact Lean-version matching & Mirror the JSON evidence
contract for LaTeX source spans now; make jixia a future optional Lean static
analysis backend \\
\addlinespace
ReProver/LeanDojo & Accessible-premise restriction, DPR-style premise
retrieval, best-first tactic search, status propagation & High for algorithm,
low for direct dependency & Build a document-level premise index restricted to
nearby/source-accessible assumptions before hypothesis expansion \\
\addlinespace
Aesop & Normalize/safe/unsafe rule phases; best-first search over goals &
High. This is a simple discipline to encode in Python & Require every expansion
provider to declare \texttt{normalize}, \texttt{safe}, or \texttt{unsafe};
unsafe operators must branch and record assumptions \\
\addlinespace
Pantograph & Machine-to-machine Lean proof-state API; MCTS-compatible search
state; goal/tactic branching & Medium after Lean environment gate & Define
\texttt{ProofStateBackend}; use Pantograph for search evidence but certify with
direct Lean \\
\addlinespace
Hammer family & Select premises, translate to ATP,
search, reconstruct proof in the assistant & High as proof discipline, not as
Python dependency & Keep external search separate from certification; every
repair must reconstruct through a deterministic backend or explicit assumption \\
\addlinespace
HTPS/Evariste & AND/OR hypergraph for tactics that create multiple subgoals &
High for data model, low for dependency & Add \texttt{AlternativeNode} and
\texttt{ObligationSetNode}; conjunctive assumption obligations must all close \\
\addlinespace
DeepSeek-Prover-V1.5/RMaxTS & Truncate-and-resume MCTS, intrinsic reward for
new proof states, root/tree parallelization & Medium for future deep profile &
Implement smoke/standard/deep budget ladder; in deep mode reward new checked
states and keep partial paths \\
\addlinespace
Thor/Baldur/DSP & Generate candidate proof, run checker, repair from verifier
feedback, accept only checked scripts & High as report-promotion policy & Add
regression tests forbidding concrete fixes without checker evidence or accepted
assumptions \\
\addlinespace
Proof-state datasets & Replayable proof-state environment and dataset
format & Medium as logging format, low as dependency & Serialize every branch,
backend attempt, source span, and blocker to replayable JSON \\
\bottomrule
\end{longtable}
\end{small}

\subsection{Most Readily Borrowable Ideas}

The highest-return borrowable ideas are not large prover models. They are
small contracts that make agents less speculative.

\paragraph{1. Summary-first retrieval.}
Borrow LeanExplore's MCP design. A MathDevMCP agent should not receive an
entire document index or full Lean source dump by default. It should call:

\begin{verbatim}
search_context_summary(query, limit, scope)
get_context_source(context_id)
get_context_assumptions(context_id)
get_context_dependencies(context_id)
get_context_location(context_id)
\end{verbatim}

This is immediately useful for LaTeX documents even before Lean integration:
the same API can search labels, equations, assumptions, definitions, theorem
statements, and nearby prose.

\paragraph{2. Accessible-premise filtering.}
Borrow ReProver's accessible-premise idea. For document repair, the accessible
premise set should include assumptions and definitions in the same theorem
environment, earlier equations in the same section, declared global
assumptions, imported notation conventions, and user-provided assumptions. It
should exclude later claims unless the document explicitly permits forward
references. This prevents the agent from proving a step using facts the text
has not established.

\paragraph{3. Decompose/retrieve/filter/judge before generation.}
Borrow LeanSearch-v2's reasoning pattern but replace free-form judging with
structured evidence. For a hard target:

\begin{enumerate}[leftmargin=*]
\item decompose the target into needed mathematical objects and subclaims;
\item retrieve candidate assumptions, definitions, lemmas, and backend routes;
\item filter by source accessibility and type/shape compatibility;
\item judge only whether the retrieved material can support the next checked
  obligation.
\end{enumerate}

Only after this should a branch synthesize a repair patch.

\paragraph{4. Source-range evidence as a first-class object.}
Borrow jixia/CoqGym's source and replay mindset. A repair proposal must carry
the exact file, label, line, source span, dependency IDs, and checker attempts.
This should become a reusable \texttt{EvidenceRef} object rather than prose in
a report.

\paragraph{5. Partial-progress tree nodes.}
Borrow DeepSeek-Prover-V1.5's truncate-and-resume lesson. If a branch derives
steps 1--3 and fails at step 4, the tree should store the successful prefix and
expand from the blocker. It should not collapse the whole attempt into
``collect more evidence.''

\subsection{Immediate Design Delta}

The next implementation plan should add three small, reusable components before
any heavyweight prover integration:

\begin{enumerate}[leftmargin=*]
\item \texttt{ContextSearchProvider}: summary-first retrieval over document
  labels, equations, assumptions, definitions, and optional Lean declarations.
\item \texttt{AccessiblePremiseIndex}: deterministic source-accessibility and
  dependency filtering for each target label or source span.
\item \texttt{EvidenceRef}: stable references for source span, retrieved
  premise, backend attempt, blocker, and accepted assumption.
\end{enumerate}

These components are the easiest borrow from the literature and code search,
and they directly attack the regression that produced hand-wavy reports.

\section{Algorithm Families}

\subsection{Sledgehammer, CoqHammer, and the Hammer Pattern}

\paragraph{Problem solved.}
Interactive proof assistants have rich libraries but users do not want to find
all premises and write all low-level proof steps manually. Hammers bridge
interactive proof assistants with external automated provers: select relevant
facts, translate a goal, run automated provers, and reconstruct a checked proof
inside the proof assistant when possible
\cite{paulson2010sledgehammer,blanchette2016hammering,czarnecki2018hammer}.

\paragraph{Algorithm sketch.}

\begin{enumerate}[leftmargin=*]
\item Input a typed proof-assistant goal and local context.
\item Select relevant premises from the library.
\item Translate the goal and premises to an ATP/SMT-friendly representation.
\item Run one or more automated provers under timeouts.
\item If a proof is found, reconstruct a proof script in the assistant.
\item Return either a checked proof/reconstruction or prover diagnostics.
\end{enumerate}

\paragraph{Relevance to \MathDevMCP.}
The hammer pattern is the right discipline: search and proposal may happen
outside the certifier, but the certifier boundary is explicit. For
\MathDevMCP{} this means the hypothesis tree can propose source-backed
assumptions and subgoals, but final proof claims must be checked by Lean or a
scoped deterministic backend.

\paragraph{Implementation lesson.}
Premise/assumption selection is a first-class component, not an afterthought.
\MathDevMCP{} should retrieve definitions, assumptions, and local equations
from the target document before expanding speculative branches.

\paragraph{Translation to document repair.}
For \MathDevMCP{}, the hammer pattern should become:

\begin{enumerate}[leftmargin=*]
\item Parse the target and collect local definitions, assumptions, labels, and
  nearby equations.
\item Select a small premise set: source-backed assumptions first, then
  deterministic route-required assumptions, then optional modeling conventions.
\item Translate the target to the strongest backend-supported form: finite
  algebra for \SymPy/\Sage{}, direct Lean statement for local formal goals, or
  typed diagnostic obligation if no backend can encode it.
\item Run multiple backend attempts under timeouts.
\item Promote only reconstructed/checkable results; otherwise return the
  exact missing premise or failed translation as a blocker.
\end{enumerate}

\subsection{Aesop: Best-First Rule Search in Lean}

\paragraph{Problem solved.}
Aesop automates ``obvious'' Lean proofs by applying a configurable set of rules.
It is useful both as general automation and as a framework for domain-specific
automation \cite{aesopgithub}.

\paragraph{Algorithm.}
Aesop builds a search tree of Lean goals. Rules can be safe or unsafe. Safe
rules are applied eagerly and are not backtracked; unsafe rules are assigned a
success probability. The tree is explored best-first: goals reached through
high-probability rules are prioritized. Before rule application, goals are
normalized with configurable normalization rules, including \texttt{simp\_all}.

\paragraph{Important mechanics.}

\begin{itemize}[leftmargin=*]
\item State: Lean goal state.
\item Action: apply a rule/tactic.
\item Safe action: deterministic progress; no backtracking.
\item Unsafe action: branch with a probability/priority.
\item Scoring: multiply parent priority by rule success probability.
\item Termination: proof found, rule space exhausted, or resource limits.
\item Debugging: trace options expose search steps and tree.
\end{itemize}

\paragraph{Code readiness.}
The Aesop repo was cloned. It is a Lean package with \texttt{lakefile.toml} and
\texttt{lean-toolchain}. It is not a Python library. Direct use requires a
matching Lean project and lake build.

\paragraph{Recommendation.}
Use Aesop as both an algorithmic template and a Lean-side tactic. The
\MathDevMCP{} hypothesis tree should copy the safe/unsafe distinction:
parser-normalization, complete source-span recovery, and dimension checking are
safe deterministic actions; adding an economic convention or general
integrability condition is unsafe and should branch.

\paragraph{Direct design requirement.}
Every \MathDevMCP{} expansion operator should declare one of three phases:
\texttt{normalize}, \texttt{safe}, or \texttt{unsafe}. A normalize operator
must not introduce a new assumption; it may only canonicalize notation,
recover a source span, normalize dimensions, or unfold definitions. A safe
operator must be provability-preserving under the current branch assumptions.
An unsafe operator may introduce a modeling convention or one of several
sufficient assumption sets, and therefore must create a branch with an
evidence-labeled hypothesis. This one rule would prevent many hand-wavy
reports: an operator that cannot justify its safety must be recorded as a
branching hypothesis, not as a fix.

\subsection{LeanDojo, ReProver, and LeanDojo-v2}

\paragraph{Problem solved.}
LeanDojo provides programmatic interaction with Lean and extraction of proof
states/tactics/premises from Lean repositories. ReProver adds
retrieval-augmented tactic generation: retrieve relevant premises and generate
tactics conditioned on the proof state and retrieved premises
\cite{yang2023leandojo,reprovergithub}. LeanDojo-v2 expands this into an
end-to-end framework with repository tracing, agents, trainers, Pantograph-based
provers, and external APIs.

\paragraph{ReProver best-first algorithm.}
The inspected \texttt{prover/search\_tree.py} and
\texttt{prover/proof\_search.py} implement:

\begin{itemize}[leftmargin=*]
\item Node statuses: \texttt{PROVED}, \texttt{FAILED}, \texttt{OPEN}.
\item Internal node state: Lean tactic state and cumulative action log
  probability.
\item Edge: a tactic from a source state to a destination state.
\item Priority: cumulative log probability.
\item Expansion: pop highest-priority open node; generate tactic suggestions;
  run each tactic in LeanDojo; create proof-finished, error, or new internal
  nodes.
\item Status propagation: a node is proved if any child is proved; failed if
  all children failed.
\item Proof extraction: recursively choose child edge with shortest
  distance-to-proof.
\item Stopping: proof found, root failed, priority queue empty, timeout, or max
  expansions.
\end{itemize}

\paragraph{What to borrow.}
This is a nearly perfect skeleton for a \MathDevMCP{} search controller, but
the state should be broader than Lean tactic state. Replace tactic generation
with expansion operators such as assumption expansion, source-span expansion,
subgoal split, backend attempt, and patch synthesis. Replace cumulative tactic
logprob with evidence-grounded branch priority.

\paragraph{Implementation skeleton to adapt.}
The \MathDevMCP{} controller can adapt the ReProver loop almost literally:

\begin{verbatim}
root = SearchNode(target=T, assumptions=A0, status="open", score=0)
queue = priority_queue([root])
while queue and not budget.exhausted():
    node = queue.pop_best()
    expansions = expansion_policy(node)
    edges = []
    for expansion in expansions:
        child = apply_expansion(node, expansion)
        child = run_available_backends(child)
        child.status = propagate_local_status(child)
        edges.append(SearchEdge(expansion, child))
        if child.status in {"open", "partial"}:
            queue.push(child)
    node.out_edges = edges
    propagate_status_to_ancestors(node)
return extract_best_branch(root), render_blockers(root)
\end{verbatim}

The differences from ReProver are the important part. The action is not a Lean
tactic; it is a repair operator. The environment is not only Lean; it is the
set of deterministic \MathDevMCP{} tools. The reward is not a model log
probability; it is evidence quality and patchability.

\paragraph{Code readiness.}
The original LeanDojo clone is marked deprecated by its README. LeanDojo-v2 is
more complete but heavy: it requires Python 3.11, PyTorch, Transformers,
DeepSpeed, TRL, PEFT, Ray, GitHub/Hugging Face tokens for common workflows, and
Pantograph for Lean RPC. Active \MathDevMCP{} Python did not import
\texttt{lean\_dojo} or \texttt{lean\_dojo\_v2} during the probe.

\paragraph{Recommendation.}
Do not make LeanDojo-v2 a required dependency of \MathDevMCP{} core. Implement
a native search controller and define an optional LeanDojo/Pantograph adapter
that can be enabled in a backend environment. Reuse the search-tree structure
and proof-attempt ledger ideas.

\subsection{PyPantograph}

\paragraph{Problem solved.}
Pantograph exposes Lean 4 as a machine-to-machine interaction environment:
start goals, execute tactics, inspect states, load sketches, check files, and
manage proof-state branching \cite{pantographgithub}.

\paragraph{Algorithm/API mechanics.}
The inspected \texttt{pantograph/server.py} exposes a \texttt{Server} with
methods such as:

\begin{itemize}[leftmargin=*]
\item \texttt{goal\_start}: create a goal state from an expression;
\item \texttt{goal\_tactic}: run a tactic on a selected goal;
\item \texttt{goal\_continue}: resume a target after branch search;
\item \texttt{expr\_type}: check expression types;
\item error handling for parse errors, tactic failure, unsafe/sorry-producing
  tactics.
\end{itemize}

The inspected \texttt{pantograph/search.py} provides:

\begin{itemize}[leftmargin=*]
\item a \texttt{SearchState} with goal state, parent, goal priorities,
  children, tried tactics, visit counts, and exhaustion flags;
\item an \texttt{Agent} interface with \texttt{next\_tactic} and
  \texttt{guidance};
\item an \texttt{MCTSAgent} interface with \texttt{estimate}, \texttt{select},
  and \texttt{backup}.
\end{itemize}

\paragraph{Code readiness.}
PyPantograph was cloned. It is small and clean, with Apache-2.0 license and a
Python package. But building requires its Lean REPL with matching Lean version:
the build script runs \texttt{lake build repl} and copies the compiled REPL into
the Python package. Active Python did not import \texttt{pantograph}; no build
was attempted.

\paragraph{Recommendation.}
This is the best direct integration candidate after an environment gate. For
the next \MathDevMCP{} phase, do not block on Pantograph. Design a
\texttt{ProofStateBackend} interface so Pantograph can later implement
\texttt{start\_goal}, \texttt{run\_tactic}, \texttt{check\_expr}, and
\texttt{branch}.

\paragraph{Concrete adapter boundary.}
The Pantograph adapter should not be allowed to certify a document repair by
itself. It should return proof-state evidence:

\begin{verbatim}
{
  "backend": "pantograph",
  "goal": "...",
  "attempted_tactics": [...],
  "result": "closed|open|tactic_failure|server_error",
  "proof_script_candidate": "...",
  "requires_direct_lean_check": true
}
\end{verbatim}

If the branch claims a proof, \MathDevMCP{} should reconstruct a standalone
Lean theorem and run direct Lean checking. Pantograph is an interactive search
environment; direct Lean is the certification boundary.

\subsection{CoqGym, HOList, GamePad, and TacticZero}

\paragraph{Problem solved.}
These systems expose theorem proving as a machine-learning environment:
state, action/tactic, transition, reward, and datasets of proof states
\cite{bansal2019holist,yang2019coqgym,huang2019gamepad,wu2021tacticzero}.

\paragraph{CoqGym mechanics.}
CoqGym extracts Coq proofs into JSON/LMDB records containing source commands,
proof steps, goals, hypotheses, environments, and tactic commands. This makes
proof search learnable and replayable. It also includes an interaction
environment and ASTactic model.

\paragraph{Code readiness.}
The CoqGym repo was cloned and is large even shallow. README explicitly says
native setup is nontrivial and recommends Docker. Native setup requires OPAM,
OCaml switch, Coq, SerAPI, CoqHammer, Ruby, LMDB, and building Coq projects.
No install was attempted.

\paragraph{Recommendation.}
Use as a data-format reference: \MathDevMCP{} should log search nodes and
backend transitions in a replayable JSON format. Do not integrate CoqGym code
for the hypothesis-search lane.

\subsection{HyperTree Proof Search and Evariste}

\paragraph{Problem solved.}
HyperTree Proof Search (HTPS) addresses neural theorem proving where actions
can create multiple subgoals. A proof search node is naturally AND/OR-like: a
tactic may produce several child goals, and the parent is solved only if all
children close \cite{lample2022htps}.

\paragraph{Algorithm.}
HTPS represents proof search as a hypergraph:

\begin{itemize}[leftmargin=*]
\item OR nodes represent goals that may be solved by one of several tactics;
\item hyperedges represent tactics;
\item each hyperedge points to a set of subgoals;
\item the search alternates expansion and value/priority estimation;
\item status propagates through the hypergraph: a hyperedge succeeds only if
  all subgoals succeed.
\end{itemize}

\paragraph{Code readiness.}
The Evariste repo was cloned. Its README explicitly says the code does not run
out of the box because Meta-internal references were removed without fixing
generic entry points. Majority license is CC-BY-NC, with some subcomponents
under separate licenses.

\paragraph{Recommendation.}
Do not integrate Evariste code. Borrow the AND/OR hypergraph idea. This matters
for \MathDevMCP{} because many derivation branches split into multiple
obligations: e.g., conditional expectation requires law, measurability, and
integrability; Bellman recursion requires state/action set, transition kernel,
reward finiteness, and terminal condition.

\paragraph{Direct design requirement.}
The search tree should not be a plain OR tree. Some expansions produce a set of
obligations that must all close. For example, the finite conditional
expectation branch succeeds only if the probability law, finiteness, and
algebraic rewrite obligations all close. This requires an AND/OR node type:
\texttt{AlternativeNode} for competing routes and \texttt{ObligationSetNode}
for conjunctive subgoals.

\subsection{Draft-Sketch-Prove, Thor, and Baldur}

\paragraph{Problem solved.}
These systems use language models to propose informal proof sketches or full
proofs, then use formal tools for checking and repair
\cite{jiang2022draft,jiang2022thor,first2023baldur}.

\paragraph{Algorithm pattern.}

\begin{enumerate}[leftmargin=*]
\item Generate a candidate proof/sketch/script.
\item Run a proof assistant or ATP-backed checker.
\item Parse failures into repair prompts or subgoals.
\item Iterate under a budget.
\item Accept only checked proof scripts; otherwise return diagnostics.
\end{enumerate}

\paragraph{Relevance.}
The important idea is not ``ask the LLM harder.'' It is the verifier loop. A
\MathDevMCP{} branch may be proposed by an agent, but every branch must be run
through parsers, \SymPy/\Sage, shape/type checks, or \Lean{} before it can be
promoted.

\subsection{Tree of Thoughts}

\paragraph{Problem solved.}
Tree of Thoughts (ToT) generalizes single-chain prompting into a search over
intermediate thought states \cite{yao2023tot}.

\paragraph{Algorithm from inspected code.}
The official \texttt{bfs.py} implements:

\begin{enumerate}[leftmargin=*]
\item Maintain candidate partial outputs \(y\).
\item Generate successors with either \texttt{sample} or \texttt{propose}.
\item Evaluate candidates with \texttt{value} or \texttt{vote} prompts.
\item Select top candidates greedily or by sampling.
\item Repeat for a fixed number of task steps.
\item Log generated candidates, values, and selected states.
\end{enumerate}

\paragraph{Code readiness.}
Repo cloned; \texttt{compileall} passed for source. It depends on the old
\texttt{openai==0.27.7} package and API keys for actual runs. Active Python did
not have \texttt{openai} or \texttt{tot} installed.

\paragraph{Recommendation.}
Do not reuse code directly. Borrow the generate/evaluate/select interface, but
replace LLM value prompts with deterministic evidence scoring.

\subsection{Language Agent Tree Search}

\paragraph{Problem solved.}
LATS combines language-model generation, environment feedback, reflection, and
tree search for reasoning/acting tasks \cite{zhou2023lats}.

\paragraph{Algorithm from inspected code.}
The programming \texttt{mcts.py} implements:

\begin{itemize}[leftmargin=*]
\item Node fields: solution, parent, children, value, visits, reflection, test
  feedback.
\item Selection: recursively choose child by \UCT.
\item Expansion: generate \(n\) new candidate solutions conditioned on prior
  solution, feedback, and reflection.
\item Simulation/evaluation: run tests, compute reward from passed tests and
  final success.
\item Backpropagation: add reward to the child and ancestors.
\item Logging: write solutions, feedback, reflection, and pass/fail status.
\end{itemize}

\paragraph{Code readiness.}
Repo cloned. It requires OpenAI/API-key paths and task-specific environments
such as HotPotQA/WebShop/programming executors. It is not proof-specific.

\paragraph{Recommendation.}
Borrow the control loop and logging discipline. Do not reuse as a dependency.
For \MathDevMCP, the environment feedback should be backend evidence, not unit
tests over generated code.

\subsection{DeepSeek-Prover, HunyuanProver, and Modern Guided Tree Search}

\paragraph{Problem solved.}
Recent prover systems combine large theorem-proving models, synthetic data,
proof-assistant feedback, reinforcement learning, and guided tree search for
\Lean{} theorem proving \cite{xin2024deepseekprover,deepseek2025v15,li2025hunyuanprover}.

\paragraph{Algorithmic implications.}
The lesson for \MathDevMCP{} is that hard proving is a high-budget search
problem. These systems rely on many generated attempts and verifier feedback.
They do not imply that a small local tool should train a prover. They do imply
that \MathDevMCP{} should record many attempted branches and should not expect
one pass to solve hard derivations.

\paragraph{Recommendation.}
Use as strategic evidence that verifier-guided search is the right direction.
Do not depend on large prover models in the core workflow.

\paragraph{Code inspection note.}
The DeepSeek-Prover-V1.5 repo exposes a high-budget RMaxTS configuration:
sample count \(6400\), concurrent sampling \(32\), Lean concurrent requests
\(64\), batch size \(512\), and \(256\) search processes. This is not a local
default for \MathDevMCP{}. It is a scale warning: genuinely hard derivations may
require many branches. The immediate implementation should therefore expose a
budget ladder:

\begin{itemize}[leftmargin=*]
\item \texttt{smoke}: 10--20 nodes, one or two deterministic routes.
\item \texttt{standard}: 80--200 nodes, multiple assumption branches and
  backend attempts.
\item \texttt{deep}: explicit user-approved run with higher node count,
  process parallelism, and optional Lean/Pantograph adapters.
\end{itemize}

\section{Decision: What Should MathDevMCP Build?}

\subsection{Do Not Pick a Single External Algorithm as Core}

No inspected codebase is a clean drop-in for \MathDevMCP{}:

\begin{itemize}[leftmargin=*]
\item LeanDojo-v2 is powerful but too heavy for a core dependency.
\item PyPantograph is promising but needs a Lean/Pantograph environment gate.
\item Aesop is Lean-side, not a Python document-repair engine.
\item Evariste explicitly does not run out of the box.
\item ToT/LATS are not proof systems and rely on LLM/API feedback.
\item CoqGym is a dataset/environment reference, not a direct integration path.
\end{itemize}

The right plan is a \MathDevMCP-native controller with optional adapters.

\subsection{Recommended Core Algorithm: Evidence-Grounded Hypothesis Tree}

\paragraph{State.}
Each node should record:

\begin{verbatim}
node_id
parent_id
depth
target
source_span
claim_kind
postulated_assumptions
accepted_assumptions
subgoals
derived_steps
backend_attempts
blockers
patch_candidates
score
status
\end{verbatim}

\paragraph{Statuses.}

\begin{itemize}[leftmargin=*]
\item \texttt{proved}: all required obligations are certified or accepted under
  explicit assumptions.
\item \texttt{refuted}: a backend/counterexample refuted the target.
\item \texttt{blocked}: expansion hit a concrete blocker.
\item \texttt{open}: branch can still be expanded.
\item \texttt{budget\_exhausted}: not solved within search budget.
\item \texttt{rejected}: malformed or unsupported candidate.
\end{itemize}

\paragraph{Actions.}

\begin{description}[style=nextline,leftmargin=0pt]
\item[Source-span expansion.]
Recover complete equation/proof span. Evidence: LaTeX parser and localization
check.
\item[Assumption expansion.]
Add finite-state, kernel-integrability, shape, Bellman, accounting, or
differentiability branch. Evidence: \texttt{assumptions\_for}, domain
templates, and source references.
\item[Subgoal split.]
Turn target into local obligations. Evidence: typed obligation IR.
\item[Algebra/backend check.]
Try \SymPy/\Sage{} on algebraic or finite-sum step. Evidence: backend result.
\item[Lean/Pantograph check.]
Try a formal Lean local theorem or tactic path. Evidence: direct Lean final
check is required before certification.
\item[Patch synthesis.]
Construct exact assumption block and replacement LaTeX. Evidence: parser and
source-span checks.
\item[Blocker record.]
Stop branch with concrete missing object. Evidence: blocker type and required
next artifact.
\end{description}

\paragraph{Scoring.}
Start with best-first/beam search, inspired by Aesop and ReProver. Use a
transparent additive score:

\begin{align*}
  S(n)={}&
  4C_{\Lean}(n)
  + 3C_{\Sage/\SymPy}(n)
  + 2C_{\mathrm{parser/type}}(n) \\
  &+ 2R_{\mathrm{source}}(n)
  + P_{\mathrm{patch}}(n)
  - 3B_{\mathrm{fatal}}(n) \\
  &- 2H_{\mathrm{unsupported}}(n)
  - M_{\mathrm{malformed}}(n).
\end{align*}

Later add \UCT:

\[
  UCT(n) = \bar r(n) + c \sqrt{\frac{\log N(parent(n))}{N(n)}}.
\]

Here reward \(\bar r(n)\) must be evidence-grounded: checked steps, fewer
missing assumptions, source-backed definitions, and patchability. It must not
be LLM confidence.

\paragraph{Expansion policy.}
For the first implementation, avoid open-ended MCTS. Use:

\begin{enumerate}[leftmargin=*]
\item deterministic root normalization;
\item generate a small fixed set of expansion operators;
\item best-first or beam search with \texttt{max\_nodes}, \texttt{max\_depth},
  \texttt{max\_branches\_per\_blocker}, and backend timeouts;
\item process-level parallelism only for side-effect-free backend attempts or
  branch scoring;
\item stable final ordering for reproducible reports.
\end{enumerate}

\subsection{Optional Adapter Interfaces}

Define adapters but keep core independent:

\begin{verbatim}
class SearchBackend:
    def can_handle(obligation) -> BackendCapability
    def attempt(node, timeout) -> BackendAttempt

class ProofStateBackend(SearchBackend):
    def start_goal(statement) -> ProofState
    def run_tactic(state, tactic) -> ProofStateResult

class ExpansionProvider:
    def expand(node, budget) -> list[Expansion]
\end{verbatim}

Concrete adapters:

\begin{itemize}[leftmargin=*]
\item \texttt{SymPyBackend}: active now.
\item \texttt{SageBackend}: active if \texttt{/usr/bin/sage} path is usable.
\item \texttt{LeanDirectBackend}: direct Lean check only; certification
  authority.
\item \texttt{PantographBackend}: optional proof-state/tactic interaction.
\item \texttt{LeanDojoV2Backend}: optional tracing/retrieval/prover backend in
  isolated environment.
\end{itemize}

\subsection{Implementation Blueprint}

\paragraph{Proposed Python entry point.}
The reusable function should be a pure-Python high-level workflow:

\begin{verbatim}
hypothesis_search_derivation(
    tex_path: str | None,
    label: str | None,
    target: str,
    context: str | None = None,
    provided_assumptions: list[str] | None = None,
    search_profile: Literal["smoke", "standard", "deep"] = "standard",
    backends: list[str] = ["sympy", "sage", "lean"],
    output_path: str | None = None,
) -> dict[str, Any]
\end{verbatim}

It should also have a Markdown renderer:

\begin{verbatim}
render_hypothesis_search_report(result) -> str
\end{verbatim}

The output path is optional. If provided, the function writes a report and
returns the same structured JSON result. This matches the successful
\texttt{audit\_math\_document\_rigor} pattern but gives the agent a stronger
within-problem search artifact.

\paragraph{Core data structures.}

\begin{verbatim}
@dataclass(frozen=True)
class SearchBudget:
    max_nodes: int
    max_depth: int
    max_branches_per_node: int
    backend_timeout_s: float
    workers: int

@dataclass
class SearchNode:
    node_id: str
    parent_id: str | None
    kind: Literal["alternative", "obligation_set", "leaf"]
    target: str
    assumptions: list[Assumption]
    obligations: list[Obligation]
    derivation_steps: list[DerivationStep]
    backend_attempts: list[BackendAttempt]
    blockers: list[Blocker]
    patch_candidates: list[PatchCandidate]
    score: float
    status: Literal[
        "proved", "refuted", "partial", "blocked",
        "open", "budget_exhausted", "rejected"
    ]
\end{verbatim}

\paragraph{Expansion providers.}

\begin{description}[leftmargin=*]
\item[Source span provider.] Takes a label or line range and returns the
  smallest complete equation/proof span. It must record file, line, section
  path, and parser confidence.
\item[Assumption provider.] Calls or generalizes \texttt{assumptions\_for} to
  produce alternative sufficient assumption sets, each with mathematical
  rationale and source/provenance.
\item[Derivation split provider.] Converts one large equality or prose claim
  into typed obligations: definition expansion, algebraic equality, shape/type
  constraint, stochastic expectation condition, or modeling convention.
\item[Backend provider.] Runs routeable obligations through \SymPy{}, \Sage{},
  direct Lean, or optional Pantograph/LeanDojo adapters.
\item[Patch provider.] Produces exact assumption blocks and replacement LaTeX
  only when all required branch steps are either checked or explicit accepted
  assumptions.
\end{description}

\paragraph{Promotion rule.}
A branch may be reported as a concrete repair only if:

\begin{enumerate}[leftmargin=*]
\item every derivation step has a checker status of \texttt{certified} or
  \texttt{accepted\_assumption};
\item every accepted assumption has a location, statement, why-needed field,
  and derivation role;
\item every backend claim has a backend attempt record;
\item every patch has the source span it replaces and the assumptions it
  introduces;
\item all remaining gaps are explicitly in the diagnostic-abstention ledger.
\end{enumerate}

\subsection{Parallelism and Search Budget}

The literature supports parallel search, but the first implementation should
use controlled process-level parallelism only where branch attempts are
side-effect-free. A good initial policy is:

\begin{itemize}[leftmargin=*]
\item expand the priority queue deterministically in the parent process;
\item submit independent backend attempts for sibling branches to a process
  pool;
\item merge results by stable node id, not by completion order;
\item cap backend timeouts separately from global search budget;
\item record every timeout as a blocker or rejected backend attempt, not as a
  mathematical failure.
\end{itemize}

This gives speed without making reports non-reproducible.

\subsection{Required Report Sections}

The Markdown/LaTeX report should include:

\begin{enumerate}[leftmargin=*]
\item \textbf{Question and location.} The target, source file, line range, and
  section path.
\item \textbf{Budget and tools.} Search profile, node/depth/backend timeouts,
  and exact backend/tool calls.
\item \textbf{Best branch.} Assumptions added, derivation steps, backend
  evidence, and patch candidates.
\item \textbf{Why assumptions are missing.} For each added assumption, show the
  failed obligation it closes.
\item \textbf{How derivation works under assumptions.} Step-by-step route with
  checked or accepted evidence status.
\item \textbf{Blocked alternatives.} Branches attempted, where they stopped,
  and the smallest artifact needed to continue.
\item \textbf{Non-claims.} No proof of global minimality, no claim about
  unsearched branches, no certification from LLM-only text, and no Lean proof
  unless direct Lean passed.
\end{enumerate}

\section{Implementation Readiness Ranking}

\begin{enumerate}[leftmargin=*]
\item \textbf{Implement now: native search controller.}
  The core can be pure Python using existing \MathDevMCP{} structures and
  deterministic backends.
\item \textbf{Implement now: deterministic expansion templates.}
  Start with conditional expectation, finite-horizon NPV, accounting identity,
  Bellman recursion, and shape/dimension constraints.
\item \textbf{Implement next: optional Pantograph adapter.}
  PyPantograph is the cleanest proof-state API but requires environment setup.
\item \textbf{Implement next: Lean direct final-check discipline.}
  Every Lean-derived branch must end in direct Lean checking.
\item \textbf{Defer: LeanDojo-v2/ReProver model integration.}
  Useful but heavy; needs separate backend env, tokens, model dependencies, and
  benchmark plan.
\item \textbf{Reference only: Evariste/HTPS code, CoqGym, ToT, LATS.}
  Borrow algorithms/data formats/control loops, not runtime dependencies.
\end{enumerate}

\section{Implementation Decision Matrix}

\begin{description}[style=nextline,leftmargin=0pt]
\item[Native evidence-grounded search controller.]
Expected value: directly fixes weak reports and is compatible with current
tools. Main risk: schema and tests must be disciplined. Decision: start here.
\item[Pantograph optional adapter.]
Expected value: Lean proof-state feedback and tactic search. Main risk:
Lean/lake/Pantograph environment gate. Decision: build after the core
controller.
\item[LeanDojo-v2 / ReProver integration.]
Expected value: retrieval and trained prover infrastructure. Main risk: heavy
dependencies, tokens, and model/runtime cost. Decision: separate backend lane
only.
\item[Aesop Lean-side tactic.]
Expected value: strong automation for generated Lean goals. Main risk: Lean
project and domain formalization required. Decision: use when a direct Lean
route exists.
\item[DeepSeek-Prover-V1.5 style search.]
Expected value: evidence for high-budget verifier-guided search. Main risk:
GPU/model/proof benchmark cost. Decision: strategic reference, not core.
\item[ToT/LATS code reuse.]
Expected value: simple tree-search control examples. Main risk: LLM/API-centric
and not proof-specific. Decision: borrow abstractions only.
\item[CoqGym/Evariste code reuse.]
Expected value: rich proof-search history. Main risk: heavy, non-running, or
mismatched stack. Decision: do not integrate.
\end{description}

The decision is not to pick a paper implementation wholesale. The decision is
to build the missing \MathDevMCP{} abstraction that all successful systems
share: explicit search state, verifier feedback, status propagation, budgeted
branching, and replayable evidence.

\section{How the Report Should Change}

A derivation-grade report should no longer say:

\begin{quote}
Derivation route: justify the reconstructed equality from context.
\end{quote}

It should say:

\begin{quote}
Search budget: 80 nodes, depth 5, 4 workers. Best branch:
finite-state conditional expectation. Source span was complete and parsed.
Step 1 expands the NPV definition; parser checked. Step 2 applies linearity of
finite conditional expectation; \SymPy{} checked the finite weighted-sum
algebra. Step 3 uses the economic zero-profit pricing convention; this remains
an explicit modeling assumption. Blocked alternatives: kernel-integrability
branch lacked a conditional kernel; Bellman branch lacked terminal boundary.
Patch: insert Assumption A and replace equation block B.
\end{quote}

This is the practical standard for implementation.

\section{Proposed Next Implementation Plan}

\begin{enumerate}[leftmargin=*]
\item Add \texttt{hypothesis\_search.py} with node/edge/budget dataclasses,
  status propagation, stable best-first queue, JSON serialization, and
  deterministic scoring.
\item Add source-span and assumption expansion providers. The first target
  templates should be finite conditional expectation, finite-horizon NPV,
  accounting identity, Bellman recursion, and shape/dimension constraints.
\item Add backend attempt provider that reuses the existing route layer for
  \SymPy/\Sage/direct Lean where available. Backend unavailable must become a
  blocker, not a refutation.
\item Add patch synthesis only after a branch satisfies the promotion rule.
  Patch records must include replacement span, inserted assumptions, and
  derivation role for each assumption.
\item Add report renderer with required sections: best branch, derivation under
  assumptions, blocked alternatives, exact tools used, budget, and non-claims.
\item Add tests that forbid concrete repairs without checked steps or accepted
  assumptions. Include regression tests for the previous hand-wavy phrases such
  as ``justify from context'' and ``collect more evidence''.
\item Integrate into \texttt{audit\_math\_document\_rigor} for selected hard
  labels behind a profile flag. Keep the old path as fallback until reports are
  stable.
\item Add optional Pantograph adapter only after the core controller passes
  tests and direct Lean readiness is reliable in a separate backend env.
\end{enumerate}

\subsection{First Experiment After Implementation}

The first application experiment should be the credit-card NPV document because
it contains exactly the kind of finite-state, accounting, discounting, and
expectation assumptions that stress the search design. The evidence contract
should be:

\begin{description}[leftmargin=*]
\item[Question.] Can the new search controller produce concrete, derivation
grade repairs instead of diagnostic-only comments for selected NPV claims?
\item[Baseline.] Current \texttt{audit\_math\_document\_rigor} report on the
same labels.
\item[Promotion criterion.] At least one target produces a best branch with
  explicit assumptions, a derivation route under those assumptions, checked or
  accepted evidence per step, and exact patch text.
\item[Veto diagnostics.] Any concrete repair lacking location, assumption role,
  backend/source evidence, or patch span is a failure.
\item[Non-claim.] Passing this experiment does not prove the whole document is
  rigorous; it only validates the search/report machinery on selected targets.
\end{description}

\section{Final Recommendation}

The implementation should not start by installing a large external prover
stack. It should start by making \MathDevMCP{} search explicit and
reproducible:

\begin{center}
\textbf{agent proposes candidate branches; deterministic tools check nodes;
the controller records evidence, blockers, scores, and budgets.}
\end{center}

This directly addresses the user's concern: hard derivation repair is a search
problem under uncertainty, not a one-pass report-generation problem.

\begin{thebibliography}{99}

\bibitem{paulson2010sledgehammer}
Lawrence C. Paulson and Jasmin C. Blanchette.
\newblock Three years of experience with Sledgehammer, a practical link between
automatic and interactive theorem provers.
\newblock In \emph{Proceedings of the 8th International Workshop on the
Implementation of Logics}, 2010.
\newblock \url{https://isabelle.in.tum.de/doc/sledgehammer.pdf}.

\bibitem{blanchette2016hammering}
Jasmin C. Blanchette, Cezary Kaliszyk, Lawrence C. Paulson, and Josef Urban.
\newblock Hammering towards QED.
\newblock \emph{Journal of Formalized Reasoning}, 9(1), 2016.

\bibitem{czarnecki2018hammer}
Lukasz Czajka and Cezary Kaliszyk.
\newblock Hammer for Coq: Automation for dependent type theory.
\newblock \emph{Journal of Automated Reasoning}, 61, 2018.

\bibitem{aesopgithub}
Jannis Limperg and Lean community contributors.
\newblock Aesop: Automated proof search for Lean.
\newblock GitHub repository.
\newblock \url{https://github.com/leanprover-community/aesop}.

\bibitem{bansal2019holist}
Kshitij Bansal, Sarah M. Loos, Markus N. Rabe, Christian Szegedy, and Stewart Wilcox.
\newblock HOList: An environment for machine learning of higher-order logic theorem proving.
\newblock In \emph{International Conference on Machine Learning}, 2019.

\bibitem{yang2019coqgym}
Kaiyu Yang and Jia Deng.
\newblock Learning to prove theorems via interacting with proof assistants.
\newblock In \emph{International Conference on Machine Learning}, 2019.

\bibitem{huang2019gamepad}
Daniel Huang, Prafulla Dhariwal, Dawn Song, and Ilya Sutskever.
\newblock GamePad: A learning environment for theorem proving.
\newblock In \emph{International Conference on Learning Representations}, 2019.

\bibitem{wu2021tacticzero}
Yuhuai Wu, Albert Q. Jiang, Wenda Li, Markus N. Rabe, Charles E. Staats,
Mateja Jamnik, and Christian Szegedy.
\newblock TacticZero: Learning to prove theorems from scratch with deep
reinforcement learning.
\newblock In \emph{NeurIPS}, 2021.

\bibitem{yang2023leandojo}
Kaiyu Yang, Aidan M. Swope, Alex Gu, Rahul Chalamala, Peiyang Song, Shixing Yu,
Saad Godil, Ryan J. Prenger, and Anima Anandkumar.
\newblock LeanDojo: Theorem proving with retrieval-augmented language models.
\newblock In \emph{Neural Information Processing Systems Datasets and
Benchmarks Track}, 2023.
\newblock \url{https://github.com/lean-dojo/LeanDojo}.

\bibitem{reprovergithub}
LeanDojo contributors.
\newblock ReProver: Retrieval-augmented theorem proving for Lean.
\newblock GitHub repository.
\newblock \url{https://github.com/lean-dojo/ReProver}.

\bibitem{pantographgithub}
Leni Aniva, Chuyue Sun, Brando Miranda, Clark Barrett, and Sanmi Koyejo.
\newblock Pantograph: A machine-to-machine interaction interface for advanced
theorem proving, high level reasoning, and data extraction in Lean 4.
\newblock 2024.
\newblock \url{https://github.com/stanford-centaur/PyPantograph}.

\bibitem{lample2022htps}
Guillaume Lample, Timothée Lacroix, Marie-Anne Lachaux, Thibaut Lavril,
Xavier Martinet, Amandine Hayat, Gabriel Ebner, Aurélien Rodriguez, and
colleagues.
\newblock HyperTree Proof Search for neural theorem proving.
\newblock In \emph{NeurIPS}, 2022.

\bibitem{jiang2022draft}
Albert Q. Jiang, Wenda Li, Jesse Michael Han, and Yuhuai Wu.
\newblock Draft, Sketch, and Prove: Guiding formal theorem provers with
informal proofs.
\newblock In \emph{International Conference on Learning Representations}, 2023.

\bibitem{jiang2022thor}
Albert Q. Jiang, Sean Welleck, Jin Peng Zhou, Timothée Lacroix, Jiacheng Liu,
Wenda Li, Mateja Jamnik, Guillaume Lample, and Yuhuai Wu.
\newblock Thor: Wielding hammers to integrate language models and automated
theorem provers.
\newblock In \emph{NeurIPS}, 2022.

\bibitem{first2023baldur}
Emily First, Markus N. Rabe, Talia Ringer, and Yuriy Brun.
\newblock Baldur: Whole-proof generation and repair with large language models.
\newblock In \emph{ACM Joint European Software Engineering Conference and
Symposium on the Foundations of Software Engineering}, 2023.

\bibitem{xin2024deepseekprover}
Huajian Xin, Daya Guo, Zhihong Shao, Zehui Ren, Qihao Zhu, Bo Liu, Chong Ruan,
Wenda Li, and Xiaodan Liang.
\newblock DeepSeek-Prover: Advancing theorem proving in LLMs through large-scale
synthetic data.
\newblock 2024.

\bibitem{deepseek2025v15}
DeepSeek-AI.
\newblock DeepSeek-Prover-V1.5: Harnessing proof assistant feedback for
reinforcement learning and Monte-Carlo tree search.
\newblock 2024/2025 technical report.

\bibitem{asher2025leanexplore}
Justin Asher.
\newblock LeanExplore: A search engine for Lean 4 declarations.
\newblock 2025.
\newblock \url{https://arxiv.org/abs/2506.11085}.
\newblock Code: \url{https://github.com/justincasher/lean-explore}.

\bibitem{gao2026leansearchv2}
Guoxiong Gao, Zeming Sun, Jiedong Jiang, Yutong Wang, Jingda Xu, Peihao Wu,
Bryan Dai, and Bin Dong.
\newblock LeanSearch v2: Global premise retrieval for Lean 4 theorem proving.
\newblock 2026.
\newblock \url{https://arxiv.org/abs/2605.13137}.
\newblock Code: \url{https://github.com/frenzymath/LeanSearch-v2}.

\bibitem{jixiagithub}
FrenzyMath contributors.
\newblock jixia: A static analysis tool for Lean 4.
\newblock GitHub repository.
\newblock \url{https://github.com/frenzymath/jixia}.

\bibitem{li2025hunyuanprover}
HunyuanProver team.
\newblock HunyuanProver: A scalable data synthesis framework and guided tree
search for automated theorem proving.
\newblock 2025 technical report.

\bibitem{yao2023tot}
Shunyu Yao, Dian Yu, Jeffrey Zhao, Izhak Shafran, Thomas L. Griffiths, Yuan Cao,
and Karthik Narasimhan.
\newblock Tree of Thoughts: Deliberate problem solving with large language
models.
\newblock In \emph{NeurIPS}, 2023.
\newblock \url{https://github.com/princeton-nlp/tree-of-thought-llm}.

\bibitem{zhou2023lats}
Andy Zhou, Kai Yan, Michal Shlapentokh-Rothman, Haohan Wang, and Yu-Xiong Wang.
\newblock Language Agent Tree Search unifies reasoning, acting, and planning in
language models.
\newblock In \emph{ICML}, 2024.
\newblock \url{https://github.com/lapisrocks/LanguageAgentTreeSearch}.

\bibitem{zheng2021minif2f}
Kunhao Zheng, Jesse Michael Han, and Stanislas Polu.
\newblock miniF2F: A cross-system benchmark for formal Olympiad-level
mathematics.
\newblock In \emph{International Conference on Learning Representations}, 2022.

\bibitem{azerbayev2023proofnet}
Zhangir Azerbayev, Bartosz Piotrowski, Hailey Schoelkopf, Edward W. Ayers,
Dragomir Radev, and Jeremy Avigad.
\newblock ProofNet: Autoformalizing and formally proving undergraduate-level
mathematics.
\newblock 2023.

\end{thebibliography}

\end{document}
